\documentclass[11pt]{article}
\usepackage[margin=1in]{geometry}
\usepackage[T1]{fontenc}
\usepackage[utf8]{inputenc}
\usepackage{amsmath,amssymb,amsthm}
\usepackage{graphicx}
\usepackage{xcolor}
\usepackage{microtype}
\usepackage[colorlinks=true,linkcolor=blue!55!black,citecolor=blue!55!black,urlcolor=blue!55!black]{hyperref}

\newtheorem{theorem}{Theorem}
\newtheorem{lemma}{Lemma}
\newtheorem{corollary}{Corollary}
\newcommand{\T}{\mathbb T}
\newcommand{\Z}{\mathbb Z}
\newcommand{\M}{\mathcal M}

\title{The conjectured sharp multiplier exponent for a flat zero patch}
\author{Candidate substantial partial-result packet for arXiv:1907.08812}
\date{11 August 2026}

\begin{document}
\maketitle

\begin{center}
\fbox{\parbox{0.92\textwidth}{\textbf{Status: candidate substantial partial result, likely valid, subject to human review.}
For $1/2<s<1$, if the Sobolev trace of $w\in H^s(\T^d)$ vanishes on a
nonempty open patch of a rational affine hyperplane, then
$1/w\notin\M_2^q$ for every $2\le q\le 1/(1-s)$.  This is exactly the
exponent conjectured by the source when the zero-set dimension is
$\sigma=d-1$.}}
\end{center}

\section{Source question and scope}

Northington studies reciprocals $u=1/w$ as $(2,q)$ Fourier multipliers and
proves that if $w\in H^s(\T^d)$ has a zero set of positive
$\sigma$-dimensional Hausdorff measure, then $u\notin\M_2^q$ in a stated
range.  The paper asks whether this range can be sharpened to
\begin{equation}\label{eq:conjectured}
 q\le \frac{d-\sigma}{d-\sigma-s}.
\end{equation}
See Section 7, Question 1, page 30 of the source \cite{source}; the page is
reproduced in Appendix~\ref{app:evidence}.

The theorem below establishes \eqref{eq:conjectured} for
$\sigma=d-1$ under the stronger geometric assumption that the zero set
contains an open flat patch in the Sobolev-trace sense.  It does not treat
arbitrary fractal zero sets of positive Hausdorff measure.

We use the source's equivalent formulation of the multiplier condition:
if $u=1/w\in\M_2^q$, then for every measurable $g$ with $wg\in L^2$,
\begin{equation}\label{eq:testing}
 \|\widehat g\|_{\ell^q(\Z^d)}\le C\|wg\|_{L^2(\T^d)}.
\end{equation}

\section*{Idea of the proof}

Near a flat zero patch, test \eqref{eq:testing} on a slab of normal thickness
$\tau$ and fixed tangential footprint.  Its Fourier transform factors, and
its $\ell^q$ norm is bounded below by $\tau^{1-1/q}$.  Since the normal
one-dimensional sections of $w$ vanish at the middle of the slab, a scaled
fractional Poincar\'e inequality bounds the other side by
$\tau^s E(\tau)$, where $E(\tau)\to0$ is the Sobolev energy in the shrinking
slab.  At $q=1/(1-s)$ the powers of $\tau$ exactly tie, and the vanishing
energy produces the endpoint contradiction.  A unimodular torus change of
variables gives every rational affine hyperplane.

\section{Two local estimates}

Write points of $\T^d$ as $(t,y)\in\T\times\T^{d-1}$.

\begin{lemma}[interval Fourier mass]\label{lem:fourier}
For $2\le q<\infty$ and $0<\tau<1/8$,
\[
 \big\|\widehat{\mathbf 1_{(-\tau,\tau)}}\big\|_{\ell^q(\Z)}
 \ge c_q\tau^{1-1/q}.
\]
\end{lemma}

\begin{proof}
For every integer $k$ with $|k|\le (8\tau)^{-1}$,
\[
 \left|\int_{-\tau}^{\tau}e^{-2\pi i kt}\,dt\right|
 =\left|\frac{\sin(2\pi k\tau)}{\pi k}\right|\ge c\tau,
\]
with the value $2\tau$ at $k=0$.  There are at least $c/\tau$ such integers,
which gives the claim after taking the $\ell^q$ norm.
\end{proof}

\begin{lemma}[localized anchored fractional Poincar\'e]\label{lem:poincare}
Let $1/2<s<1$.  If $f\in H^s(-2\tau,2\tau)$ has its continuous
representative satisfying $f(0)=0$, then
\[
 \int_{-\tau}^{\tau}|f(t)|^2\,dt
 \le C_s\tau^{2s}
 \int_{-2\tau}^{2\tau}\int_{-2\tau}^{2\tau}
 \frac{|f(t)-f(r)|^2}{|t-r|^{1+2s}}\,dr\,dt.
\]
\end{lemma}

\begin{proof}
After scaling, this is the standard one-dimensional fractional Sobolev
estimate
\[
 |F(x)-F(0)|\le C_s |x|^{s-1/2}[F]_{H^s(-2,2)},
 \qquad |x|\le1,
\]
integrated over $x\in(-1,1)$.  The point value is legitimate because
$s>1/2$.
\end{proof}

\section{Sharp exponent on a flat patch}

\begin{theorem}\label{thm:coordinate}
Let $d\ge2$, $1/2<s<1$, and $w\in H^s(\T^d)$.  Suppose there is a
nonempty open set $U\subset\T^{d-1}$ such that the Sobolev trace of $w$ on
$\{0\}\times\T^{d-1}$ vanishes almost everywhere on $U$.  Then
\[
 \frac1w\notin\M_2^q(\T^d)
 \qquad\text{for every}\qquad
 2\le q\le\frac1{1-s}.
\]
\end{theorem}

\begin{proof}
Choose a nonzero $\phi\in C_c^\infty(U)$.  The Fourier characterization of
$H^s(\T^d)$ and Fubini give
\begin{equation}\label{eq:sectionenergy}
 \int_{\T^{d-1}}[w(\cdot,y)]_{H^s(\T)}^2\,dy
 \asymp \sum_{(k_1,k')\in\Z^d}|k_1|^{2s}|\widehat w(k_1,k')|^2
 \le \|w\|_{H^s(\T^d)}^2.
\end{equation}
Thus almost every normal section belongs to $H^s(\T)$, its value at zero
equals the Sobolev trace, and Lemma~\ref{lem:poincare} applies on the support
of $\phi$.

For small $\tau>0$, put
\[
 g_\tau(t,y)=\mathbf1_{(-\tau,\tau)}(t)\phi(y).
\]
The Fourier transform factors.  Lemma~\ref{lem:fourier} therefore gives
\begin{equation}\label{eq:lower}
 \|\widehat g_\tau\|_{\ell^q(\Z^d)}
 =\|\widehat{\mathbf1_{(-\tau,\tau)}}\|_{\ell^q(\Z)}
   \|\widehat\phi\|_{\ell^q(\Z^{d-1})}
 \ge c_{q,\phi}\tau^{1-1/q}.
\end{equation}

Define the local sectional energy
\[
 E(\tau)^2=
 \int_{\operatorname{supp}\phi}
 \int_{-2\tau}^{2\tau}\int_{-2\tau}^{2\tau}
 \frac{|w(t,y)-w(r,y)|^2}{|t-r|^{1+2s}}\,dr\,dt\,dy.
\]
By \eqref{eq:sectionenergy}, the full integrand is integrable.  Absolute
continuity of the integral gives $E(\tau)\to0$ as $\tau\downarrow0$.
Integrating Lemma~\ref{lem:poincare} in $y$ yields
\begin{equation}\label{eq:upper}
 \|wg_\tau\|_2\le C_{s,\phi}\tau^sE(\tau).
\end{equation}

Assume, for contradiction, that $1/w\in\M_2^q$; in particular $w$ is
nonzero almost everywhere, so the source's testing identity applies without
ambiguity.  Apply
\eqref{eq:testing} to $g_\tau$ and combine \eqref{eq:lower}--\eqref{eq:upper}:
\[
 \tau^{1-1/q-s}\le C E(\tau).
\]
If $q<1/(1-s)$, the exponent on the left is negative, so the left side
diverges.  At the critical value $q=1/(1-s)$, the left side is one while
$E(\tau)\to0$.  Both cases are impossible.
\end{proof}

\begin{corollary}[rational affine hyperplanes]\label{cor:rational}
Let $a\in\Z^d$ be primitive and let
$H_{a,c}=\{x\in\T^d:a\cdot x=c\pmod1\}$.  If the Sobolev trace of
$w\in H^s(\T^d)$ vanishes on a nonempty relatively open patch of
$H_{a,c}$, then the conclusion of Theorem~\ref{thm:coordinate} holds.
\end{corollary}

\begin{proof}
Extend $a$ to a unimodular integer matrix.  The resulting torus automorphism,
followed by a translation, sends $H_{a,c}$ to a coordinate hyperplane.
Composition preserves Lebesgue norms, permutes Fourier coefficients by a
lattice bijection, preserves the $(2,q)$ multiplier property, and acts
boundedly and invertibly on $H^s$.  Apply Theorem~\ref{thm:coordinate}.
\end{proof}

\section{Verification, limitations, and novelty}

The proof was audited at the two potentially sharp points.  First, the slab
Fourier lower bound counts $\asymp\tau^{-1}$ coefficients of size
$\asymp\tau$.  Second, the endpoint does not follow from a global Poincar\'e
constant: it specifically uses absolute continuity of the \emph{localized}
sectional energy, which tends to zero.  The assumption $s>1/2$ is exactly what
makes the normal trace and anchored one-dimensional inequality available.

This packet proves the conjectured exponent only for a flat trace-zero patch.
It does not settle arbitrary sets with
$\mathcal H^{d-1}(\Sigma(w))>0$, lower-dimensional or fractal zero sets, the
$s=1$ boundary, or the other open problems in Section 7.

The bounded novelty check on 11 August 2026 searched the run registry and all
four cheap indexes for the arXiv id, title, sharp exponent, hyperplane, and
zero-set phrases.  Exact-title and close-phrase primary-source web searches
found the source and unrelated multiplier papers, but no later resolution or
this flat-patch theorem.  This is not an exhaustive priority claim.

\textbf{Human-review recommendation.} Check the localized anchored fractional
Poincar\'e estimate and the identification of the sectional value at $t=0$
with the Sobolev trace.  These are standard for $s>1/2$ and are the only
functional-analytic inputs beyond the source's multiplier testing identity.

\appendix
\section{Source-page evidence}\label{app:evidence}

\begin{figure}[p]
\centering
\includegraphics[height=0.80\textheight]{figures/open_problem_crop.png}
\caption{Page 30 of arXiv:1907.08812.  Question 1 asks whether the range can
be sharpened to $q\le(d-\sigma)/(d-\sigma-s)$.}
\end{figure}

\begin{thebibliography}{9}
\bibitem{source}
M. Northington,
\emph{Uncertainty Principles for Fourier Multipliers},
arXiv:1907.08812 [math.FA], 2019.
\end{thebibliography}

\end{document}
